

Large pre-trained vision models are routinely adapted to new domains/applications using lightweight adapters. Existing mechanistic interpretability analyses reuse a \emph{generic} Sparse Autoencoder (SAE) pretrained using a large dataset $\cal D$ to explain adapted variants. In this paper, we analyze the changes induced by the adapters on mechanistic interpretability. We show that the interpretability of the generic SAE degrades as the target domain moves away from $\cal D$, using multiple encoders (CLIP, SigLIP, DINO, ALIGN), adapters (LoRA, MaPLe, TAP, VPT), and SAEs (Gated, Top-K, JumpReLU, Matryoshka) using several metrics. The generic SAE features lose both reconstruction fidelity and concept selectivity as the distribution distance of the target shifts from ImageNet-like benchmarks to far ones like PathMNIST, ChestMNIST, and EuroSAT. We also show that {\em finetuned} SAEs built on the adapted activations using the target dataset recover and surpass the interpretability profile, yielding sparser and more monosemantic features. We introduce a new \emph{Domain-Aware Monosemanticity Score} (DAMS) to measure concept coherence relative to the target distribution, providing a method to determine when reuse of generic SAEs is justified and when finetuned SAEs are necessary. Going beyond correlational evidence, we intervene directly on SAE latents and find that generic-SAE latents have no detectable causal effect on the adapted model's predictions, while finetuned-SAE latents do, with the gap in how far a dictionary's latents are \emph{necessary} to correct classification tracking domain distance almost perfectly. We further show that embedding-space distance is the wrong quantity to predict degradation from --- a text-space measure of how far a target vocabulary falls outside ImageNet's is a substantially better predictor --- and combine the two into a rule that calls retrain-versus-reuse correctly in seven of eight domains out-of-sample, using only quantities measurable before any SAE is trained. Our study finds that adapter families differ mechanistically: prompt-based methods primarily steer existing features, while weight-space adaptation reshapes the feature dictionary itself. We conclude that trustworthy mechanistic interpretation of adapted models requires domain-matched SAEs as the target domain moves farther.

%Vision-Language Models (VLMs) such as CLIP, SigLIP, and ALIGN, along with self-supervised encoders like DINOv2, are routinely adapted to new domains using lightweight parameter-efficient techniques. While these adapters deliver strong empirical gains, the internal changes they induce remain poorly understood, and existing mechanistic interpretability analyses overwhelmingly rely on a single \emph{generic} Sparse Autoencoder (SAE) trained on the unadapted model and reused to explain every adapted variant, an assumption never stress-tested. We show across four encoders, four adapter families, and multiple SAE variants that generic SAEs do not transfer faithfully: their interpretability degrades substantially under adaptation, with the gap growing non-linearly in the distance between pretraining and target domains, a phenomenon we quantify using a new \emph{Domain-Aware Monosemanticity Score} (DAMS) that measures concept coherence relative to the target distribution. On datasets ranging from ImageNet-like benchmarks to far-OOD medical (MedMNIST, ChestMNIST) and satellite (EuroSAT) imagery, generic-SAE features lose both reconstruction fidelity and concept selectivity when applied to adapted CLIP, SigLIP, ALIGN, and DINOv2. We show that \emph{Fine-tuned} SAEs, trained on the adapted model, recover and surpass the original interpretability profile, yielding sparser and more monosemantic features whose gain over generic SAEs scales with domain distance. We further find that adapter families differ mechanistically: prompt-based methods primarily steer existing features, while weight-space adaptation more aggressively reshapes the feature dictionary itself. Our findings underscore the brittleness of current SAE-based interpretation pipelines and argue that trustworthy mechanistic interpretation of adapted foundation models requires domain-matched SAEs, with DAMS providing a principled criterion for when generic-SAE reuse is justified.
